
\documentclass[a4paper]{book}

%Пакеты для математических символов:
\usepackage{amsmath} % американское математическое сообщество.
\usepackage{amssymb} % миллион разных значков и готический, ажурный шрифты.
\usepackage{amscd} % диаграммы, графики.
\usepackage{amsthm} % окружения теорем, определений и тд.
\usepackage{physics} % основные физические символы
%\usepackage{latexsym} % треугольники и пьяная стрелка.

%пакеты для шрифтов:
%\usepackage{euscript} % прописной шрифт с завитушками.
\usepackage{MnSymbol} % Значеки доказательства
\usepackage{verbatim} % улучшенный шрифт "пишущей машинки".
\usepackage{array} % более удобные таблицы.
%\usepackage{multirow} % мультистолбцы в таблицах.
%\usepackage{longtable} % таблицы на несколько страниц.
%\usepackage{latexsym}

\usepackage{etoolbox}
\usepackage{slashbox} %Разделениени текста \backslashbox{}{}
\usepackage{collectbox} % Добавляет коробочки, можно складывать туда текст)

\usepackage[a4paper]{geometry}
\geometry{left=2cm}
\geometry{right=2cm}

\usepackage{hyperref} % Ссылки как внешние так и внутренние
\hypersetup{
    colorlinks=true,
    linkcolor=black,
    filecolor=magenta,      
    urlcolor=cyan,
    pdftitle={Overleaf Example},
    pdfpagemode=FullScreen,
    }
    
%Пакеты для оформления:
\RequirePackage[center, medium]{titlesec}% Стиль секций и заголовков
%\usepackage[x11names]{xcolor} % 317 новых цветов для текста.
\usepackage{float} % Позволяет использовать H, h! для локации фигур
%\usepackage{multicol} % набор текста в несколько колонн.
\usepackage{graphicx} % расширенные возможности вставки стандартных картинок.
\usepackage{subcaption} % возможность вставлять картинки в строчку
%\usepackage{caption} % возможность подавить нумерацию у caption.
\usepackage{wrapfig} % вставка картинок и таблиц, обтекаемых текстом.
\usepackage{cancel} % значки для сокращения дробей, упрощения, стремления.
\usepackage{misccorr} % в заголовках появляется точка, но при ссылке на них ее нет.
%\usepackage{indentfirst} % отступ у первой строки раздела
%\usepackage{showkeys} % показывает label формул над их номером.
%\usepackage{fancyhdr} % удобное создание верхних и нижних колонтитулов.
%\usepackage{titlesec} % еще одно создание верхних и нижних колонтитулов
\usepackage{hyperref} % Ссылки как внешние так и внутренние
\hypersetup{
    colorlinks=true,
    linkcolor=black,
    filecolor=magenta,      
    urlcolor=cyan,
    pdftitle={Overleaf Example},
    pdfpagemode=FullScreen,
    }
\usepackage{xcolor} %Позволяет перекрасить все страници
\definecolor{mycolor}{RGB}{244,228,215} %Цвет перекраски


%Пакеты шрифтов, кодировок. НЕ МЕНЯТЬ РАСПОЛОЖЕНИЕ.
\usepackage[utf8]{inputenc} % кодировка символов.
%\usepackage{mathtext} % позволяет использовать русские буквы в формулах. НЕСОВМЕСТИМО С tempora.
\usepackage[T1, T2A]{fontenc} % кодировка шрифта.
\usepackage[english, russian]{babel} % доступные языки.

%tikz:
\usepackage{tikz} % tikz
\usetikzlibrary{decorations.text} % позволяет делать текст вдоль кривой.
\usetikzlibrary{external} % позволяет кэшировать рисунки tikz.


\usepackage{pgfplotstable}

%Отступы и поля:
%размеры страницы А4 11.7x8.3in
\textwidth=7.3in % ширина текста
\textheight=10in % высота текста
\oddsidemargin=-0.5in % левый отступ(базовый 1дюйм + значение)
\topmargin=-0.5in % отступ сверху до колонтитула(базовый 1дюйм + значение)


%Сокращения
%Скобочки
\newcommand{\inrad}[1]{\left( #1 \right)}
\newcommand{\inner}[1]{\left( #1 \right)}
\newcommand{\infig}[1]{\left\{ #1 \right\}}
\newcommand{\insqr}[1]{\left[ #1 \right]}
\newcommand{\ave}[1]{\left\langle #1 \right\rangle}



%% Красивые <= и >=
\renewcommand{\geq}{\geqslant}
\renewcommand{\leq}{\leqslant}

%%Значек выполнятся
\newcommand{\per}{\hookrightarrow}

%%Векторная алгебра
\newcommand{\rot}{\text{rot}}
\renewcommand{\div}{\text{div}}
\renewcommand{\grad}{\text{grad}}

%% Более привычные греческие буквы
\renewcommand{\phi}{\varphi}
\renewcommand{\epsilon}{\varepsilon}
\newcommand{\eps}{\varepsilon}
\newcommand{\com}{\mathbb{C}}
\newcommand{\re}{\mathbb{R}}
\newcommand{\nat}{\mathbb{N}}
\newcommand{\stp}{$\filledmedtriangleleft$}
\newcommand{\enp}{$\filledmedsquare$}

%%Тензорный анализ ОТО теория поля
\newcommand{\Li}[1]{\mathfrak{L}_{#1}}
\newcommand{\crist}[3]{\cfrac{1}{2} \inner{g_{#1#2,#3} + g_{#1#3,#2} - g_{#2#3,#1}}}
\newcommand{\piv}[2]{\cfrac{\partial #1}{\partial #2}}

\makeatletter
\newcommand{\sqbox}{%
    \collectbox{%
        \@tempdima=\dimexpr\width-\totalheight\relax
        \ifdim\@tempdima<\z@
            \fbox{\hbox{\hspace{-.5\@tempdima}\BOXCONTENT\hspace{-.5\@tempdima}}}%
        \else
            \ht\collectedbox=\dimexpr\ht\collectedbox+.5\@tempdima\relax
            \dp\collectedbox=\dimexpr\dp\collectedbox+.5\@tempdima\relax
            \fbox{\BOXCONTENT}%
        \fi
    }%
}
\makeatother
\newcommand{\mergelines}[2]{
\begin{tabular}{llp{.5\textwidth}}
#1 \\ #2
\end{tabular}
}
\newcommand\tab[1][0.51cm]{\hspace*{#1}}
\newcommand\difh[2]{\frac{\partial #1}{\partial #2}}
\newcommand{\messageforpeople}[1]{HSE Faculty of Physics \ \ HSE Faculty of Physics HSE Faculty of Physics \ \ HSE Faculty of Physics HSE Faculty of Physics \ \ HSE Faculty of Physics HSE Faculty of Physics \ \ HSE Faculty of Physics HSE Faculty of Physics \ \ HSE Faculty of Physics HSE Faculty of Physics \ \ HSE Faculty of Physics HSE Faculty of Physics \ \ HSE Faculty of Physics HSE Faculty of Physics \ \ HSE Faculty of Physics }


\numberwithin{equation}{section}

\begin{document}


% \pagecolor{mycolor}
%\pagestyle{empty}
\oddsidemargin=-1.18in
\topmargin=-1.5in

\begin{titlepage}
\begin{tikzpicture}
    \draw [white] (0, 0) ellipse (4.1in and 5.7in);
    \path [postaction={
    decoration={text along path,
    text={\messageforpeople{}},
    text align=center,
    raise=0mm,
    reverse path}, decorate}]
    [postaction={
    decoration={text along path,
    text={\messageforpeople{}},
    text align=center,
    raise=3mm,
    reverse path}, decorate}]
    [postaction={
    decoration={text along path,
    text={\messageforpeople{}},
    text align=center,
    raise=6mm,
    reverse path}, decorate}]
    [postaction={
    decoration={text along path,
    text={\messageforpeople{}},
    text align=center,
    raise=9mm,
    reverse path}, decorate}]
    [postaction={
    decoration={text along path,
    text={\messageforpeople{}},
    text align=center,
    raise=12mm,
    reverse path}, decorate}] 
    [postaction={
    decoration={text along path,
    text={\messageforpeople{}},
    text align=center,
    raise=15mm,
    reverse path}, decorate}]
    [postaction={
    decoration={text along path,
    text={\messageforpeople{}},
    text align=center,
    raise=18mm,
    reverse path}, decorate}]
    [postaction={
    decoration={text along path,
    text={\messageforpeople{}},
    text align=center,
    raise=21mm,
    reverse path}, decorate}] 
    [postaction={
    decoration={text along path,
    text={\messageforpeople{}},
    text align=center,
    raise=24mm,
    reverse path}, decorate}]
    [postaction={
    decoration={text along path,
    text={\messageforpeople{}},
    text align=center,
    raise=27mm,
    reverse path}, decorate}]
    [postaction={
    decoration={text along path,
    text={\messageforpeople{}},
    text align=center,
    raise=30mm,
    reverse path}, decorate}] 
    [postaction={
    decoration={text along path,
    text={\messageforpeople{}},
    text align=center,
    raise=33mm,
    reverse path}, decorate}]
    [postaction={
    decoration={text along path,
    text={\messageforpeople{}},
    text align=center,
    raise=36mm,
    reverse path}, decorate}]
    [postaction={
    decoration={text along path,
    text={\messageforpeople{}},
    text align=center,
    raise=39mm,
    reverse path}, decorate}] 
    [postaction={
    decoration={text along path,
    text={\messageforpeople{}},
    text align=center,
    raise=42mm,
    reverse path}, decorate}]
    [postaction={
    decoration={text along path,
    text={\messageforpeople{}},
    text align=center,
    raise=45mm,
    reverse path}, decorate}]
    [postaction={
    decoration={text along path,
    text={\messageforpeople{}},
    text align=center,
    raise=48mm,
    reverse path}, decorate}] 
    [postaction={
    decoration={text along path,
    text={\messageforpeople{}},
    text align=center,
    raise=51mm,
    reverse path}, decorate}]
    [postaction={
    decoration={text along path,
    text={\messageforpeople{}},
    text align=center,
    raise=54mm,
    reverse path}, decorate}]
    [postaction={
    decoration={text along path,
    text={\messageforpeople{}},
    text align=center,
    raise=57mm,
    reverse path}, decorate}] 
    [postaction={
    decoration={text along path,
    text={\messageforpeople{}},
    text align=center,
    raise=60mm,
    reverse path}, decorate}]
    [postaction={
    decoration={text along path,
    text={\messageforpeople{}},
    text align=center,
    raise=63mm,
    reverse path}, decorate}]
    [postaction={
    decoration={text along path,
    text={\messageforpeople{}},
    text align=center,
    raise=66mm,
    reverse path}, decorate}] 
    [postaction={
    decoration={text along path,
    text={\messageforpeople{}},
    text align=center,
    raise=69mm,
    reverse path}, decorate}]
    [postaction={
    decoration={text along path,
    text={\messageforpeople{}},
    text align=center,
    raise=72mm,
    reverse path}, decorate}]
    (0, 0) ellipse (2.75in and 4.96in);
    \node at (0, 3.2in) {\Large\textbf{Домашняя работа 1}};
    \node at (0, 0in) {\includegraphics[height=4in]{FOPF.png}};
    \node at (0, -3in) {Авторы:};
    \node at (0, -3.2in) {Карибджанов Матвей};
    \node at (0, -3.5in) {2023};
\end{tikzpicture}

%2.75in and 4.96in
\end{titlepage}

\oddsidemargin=-0.5in
\topmargin=-0.5in
\pagestyle{plain}

\newpage
\tableofcontents
\newpagestyle{main}{
\setfootrule{0.4pt}
\setfoot{}{\thepage}{\chaptertitle }}
\pagestyle{main}



\chapter{Ведение в математику p-адических чисел}
\section{Строение пространства}

\subsection*{Кольцо}

Любое p-адичесое число $x$ можно предстваить как: 
\begin{equation}
    \label{def:p_adic}
    x = p^\gamma \cfrac{m}{n}
\end{equation}
Где $p \in \mathbb{N}, \ n \in \mathbb{N}, 
\ m \in \mathbb{Z}: \nexists y \in \mathbb{N} 
\per \cfrac{p}{y} \in  \mathbb{N}$, 
а также $p, \ n, \ m$ взаимно просты. 
Любое p-адическое чиcло определено толко в пространстве $\mathbb{Q}_p$,
где $p$ тоже что использовалось в определениие \ref{def:p_adic}. 
Любое число из $\mathbb{Q}_p$ можно записать, как число из $\mathbb{Q}$
и тогда сложение и произведение будет определно также как в $\mathbb{Q}$.
Конечно же такие операции не выводят нас из кольца $\mathbb{Q}_p$:

\begin{eqnarray}
    \label{eq:prod}
    p^{\gamma_1} \cfrac{m_1}{n_1} \cdot p^{\gamma_2} \cfrac{m_2}{n_2} = p^{\gamma_2+\gamma_1} \inner{\cfrac{m_2\cdot m_1}{n_2 \cdot n_1}}\\
    \label{eq:add}
    p^{\gamma_1} \cfrac{m_1}{n_1} + p^{\gamma_2} \cfrac{m_2}{n_2} = p^{\gamma_2} \inner{p^{\gamma_1} \cfrac{m_1}{n_1} + \cfrac{m_2}{n_2}}
\end{eqnarray}

Для \ref{eq:prod} так как $p$ - простое то мы не сможем выделить 
из $\inner{\cfrac{m_2\cdot m_1}{n_2 \cdot n_1}}$ еще $p$.
Для \ref{eq:add} я пологал что $\gamma_1 < \gamma_2$, 
поэтому по анологии я уже не смогу вынести $p$ из скобок.

\subsection*{Норма и метрика}

Для построения протсранства в первую очередь требутся норма: 

\begin{equation}
    \label{def:norm}
    \abs{x}_p = p^{-\gamma}
\end{equation}

Очень легко проверятся, что для такой 
нормы выполняются все необходимы критерии, используя уже доказанные мною 
првила сложения и умножения \ref{eq:prod} и \ref{eq:add}:
\begin{itemize}
    \item[1] $\abs{x}_p \geq 0$
    \item[2] $\abs{x+y}_p \leq \abs{x}_p + \abs{y}_p$
    \item[3] $\abs{xy}_p = \abs{x}_p \abs{y}_p$
\end{itemize} 

Пример:

\begin{eqnarray*}
    \abs{\cfrac{1}{2}}_2 = 2\\
    \abs{6}_2 = \cfrac{1}{2}\\
    \abs{13}_4 = 1
\end{eqnarray*}

Теперь когда мы задали норму в протсранстве $\mathbb{Q}_p$. 
Таким образом любое чило из $\mathbb{Q}$ можно записать в следующем виде:

\begin{equation}
    \label{def:sum}
    x = \sum_{\gamma}^{\infty} x_i p^i
\end{equation}

Где $x_i$ число в $p$ ситеме исчисления, то есть $x_i \in \mathbb{N},\ x_i < p$. 
Может возникнуть вопрос, как такой ряд может сойтись, здесь надо не забывать, что 
ма работаем в норме $\abs{x}_p$ тоесть ряд будет иметь норму $p^{-\gamma}$.
Приэтом чиcла в таком виде записыватся следющим образом:

\begin{equation}
    x = \overline{x_{\infty} ... x_{\gamma + 1} x_{\gamma} ... x_{0}. x_{-1} ... x_{-\gamma}}_p
\end{equation}

Например:

\begin{eqnarray*}
    123 = 123_{10}\\
    123 = 11120_3\\
    -123 = 9...999877_{10} = (9)877_{10}\\
    -\cfrac{1}{49} = (6).66_7
\end{eqnarray*}

\begin{equation}
    \label{def:metric}
    \rho_p\inner{x, y} = \abs{x - y}_p  
\end{equation}

Все критерии мерики очевидно выполнятся для $\rho_p$. Напомню
что это:
\begin{itemize}
    \item[1] $\rho\inner{x, y} \geq 0, \   \rho\inner{x, x} = 0$
    \item[2] $\rho\inner{x, y} = \rho\inner{y, x}$
    \item[3] $\rho\inner{x, y} \leq \rho\inner{x, z} + \rho\inner{z, y}$ 
\end{itemize} 


Но у данной метрики есть одно важное свойство, котрое позже пригодится 
нам при расчете интегралов:

\begin{equation}
    \label{def:ultra}
    \rho \inner{x, y} \leq max\inner{\rho \inner{x, z}, \rho \inner{z, y}}
\end{equation}

По сути я доказал это в \ref{def:sum}. Пространства 
чья метрика удовлетворят условию \ref{def:ultra} называтся 
ультраметрическими.



\section{Предельный переход $\mathbb{Q}_p \to \mathbb{Q}$}

Для понимания работы $p$ - адических чисел может быть очень полезно понимать,
что $\mathbb{Q}_p$ есть более общий случай $\mathbb{Q}$.

Можно лекко заметить что $\mathbb{Q} = \mathbb{Q}_{\infty}$, действительно для 
$\forall x \in \mathbb{Q}$ при $p = \infty, \nexists \gamma \neq 0$ 
чтобы предствить число в виде \ref{def:p_adic}. Тогда $\gamma = 0$ 
получается из \ref{def:p_adic} мы получим:

\begin{equation}
    p^0 \cfrac{m}{n} = \cfrac{m}{n}
\end{equation}

Что совпадает с определением $\mathbb{Q}$. И так выполнятся равентво

\begin{equation}
    \prod_2^{\infty} \abs{x}_i = 1
\end{equation}

Где $\abs{x}_{\infty} = \abs{x}$ тоесть норме р $\mathbb{Q}$ соответственно.

\section{Шары в $\mathbb{Q}_p$}


Шар в $\mathbb{Q}_p$ задается как это принято в мат. анлизе:

\begin{equation}
    \label{def:sp}
    B_\gamma \inner{a} = \infig{x \in \mathbb{Q}_p  \ \vline \ \rho \inner{x, a} \leq p^\gamma}
\end{equation}

Шары в $\mathbb{Q}_p$ обладают очнень не привычными нам примитивным обитателям из $\mathbb{R}_{3, 1}$ свойствами:
\begin{itemize}
    \item[1] Любая тока шара являтся ее центром
    \item[2] $\forall B_1 , B_2; B_1 \cap B_2 \neq 0 \per B_1 \in B_2 \ \lor B_2 \ \in B_1$
    \item[3] $\forall B_1$ открыт
\end{itemize}

Первое свойство это следствие ультра метричности пространства:
\begin{equation*}
    \forall x \in B_\gamma \inner{a} \per r = \rho \inner{a, z} 
    \leq max\inner{\rho \inner{x, a}, \rho \inner{x, z}}  
\end{equation*}

Так:

\begin{equation*}
    \rho \inner{x, a} \leq max \inner{\abs{z}_p, \abs{a}_p}, \rho \inner{x, z} \leq max \inner{\abs{x}_p, \abs{z}_p} \Rightarrow \rho \inner{x, a} \geq \rho \inner{x, z}
\end{equation*}

\begin{equation*}
    r = \rho \inner{a, z} \leq \rho \inner{x, a}, \ \rho \inner{a, z} \geq \rho \inner{x, a} \Rightarrow \rho \inner{a, z} = \rho \inner{x, a}
\end{equation*}

Второе свойство, можно легко доказать используя первое свойство:
\begin{equation*}
    \forall a \in B_1 \cap B_2,\forall z_1 \in B_1,\forall z_2 \in B_2 \Rightarrow B_1 := B_{r_1} \inner{a}, B_2 := B_{r_2} \inner{a}
\end{equation*}

Тоесть в качестве центра каждого шара мы можем выбрать любую точку из пересечения, допуская что $r_1 \leq r_2$

\begin{equation*}
    \rho \inner{a, z_1} \leq r_1 \leq r_2, \rho \inner{a, z_2} \leq r_2 \Rightarrow z_1 \in B_1, B_2
\end{equation*}

Третье свойство, я не буду доказывать, но отмечу, что так как кольцо $\mathbb{Q}$ не содержит 
своих предельных точек следовательно $\mathbb{Q}_p$ не содержит своих передльных точект, значит подможетво 
$B$ тоже не будет содержать предельных точек, что как мы знаем соответствует открытости множества.

\section{Интегрирование}


\subsection*{Мера Хаара}












\end{document}